\begin{appendices}

% ──────────────────────────────────────────────────────────────────────────────────────────────────────────
\chapter{Self-Appraisal}
\label{app:selfappraisal}

% ──────────────────────────────────────────────────────────────────────────────────────────────────────────
\section{Critical Self-Evaluation}

All four project objectives were met. The L-system interpreter supports all seven grammar variants; eight
distinct plant grammars are implemented; environmental response covers phototropism, gravitropism, wind, and
seasons; and the interactive visualization maintains interactive frame rates across all scenes. The ecosystem
placement system, which was not in the original objectives, was added as an extension and represents the most
technically novel contribution of the project.

\medskip
The most rigorously validated component is the L-system interpreter. The unit test suite covers all rule
classes, including edge cases for context-sensitive priority ordering and stochastic probability boundaries,
providing high confidence in the correctness of the rewriting engine. By contrast, the plant grammars
themselves were validated only visually: while the grammars for \textit{Capsella}, \textit{Lychnis}, and
\textit{Mycelis} closely match the figures in \textit{The Algorithmic Beauty of
Plants}~\cite{prusinkiewiczAlgorithmicBeautyPlants1990}, no quantitative comparison of branch angles or
internode lengths against measured botanical data was performed. This limits the claim of biological accuracy
to visual plausibility.

\medskip
The environmental response system achieves the stated objectives of interactive, per-species seasonal
behaviour, but is simplified relative to biological reality. Tropism is applied as a post-processing
deflection of the turtle frame rather than as feedback into the rewriting process, meaning individual branches
cannot respond to local light conditions. Season-driven dormancy is encoded as a simple age-scaling function
rather than being derived from accumulated temperature or photoperiod, which would be more biologically
accurate but significantly more complex to implement. These simplifications were deliberate design choices
given the time constraints and complexity of real world ecosystems however further development would allow
improved realism and the ability to visualize more complex interactions.

\medskip
The rendering pipeline fulfils the interactive frame rate requirement but lacks features that would improve
visual realism: specifically, shadow casting, antialiasing, and support for more textures and more varied
meshes. Both were identified early as lower-priority concerns relative to botanical correctness and
ecological placement, and deferred accordingly. Their absence is the most visible gap between the project's
output and commercial-quality vegetation renderers.

% ──────────────────────────────────────────────────────────────────────────────────────────────────────────
\section{Personal Reflection and Lessons Learned}

The most significant management challenge was scope creep. The original four objectives described an L-system
interpreter, a minimum of three plant grammars, environmental response, and interactive visualization. All of those
objectives excluding the interpreter expanded beyond the initial estimate: eight grammars were developed
going far beyond the required three, environmental response grew from phototropism, gravitrophism and space
constraints to include a full seasonal cycle with per-species phenology, and adaptive LOD and frustum culling
were implemented to permit the visualization of more complex scenes. Beyond the core objectives, the
ecosystem placement system and the performance logger were also unplanned additions. Individually, each
addition was small and improved the project in a clear and immediately understandable way; collectively they
consumed a substantial fraction of available development time.

\medskip
The lesson is that every feature added to a simulation of natural systems tends to reveal further features it
depends on or that would enhance it: a season system immediately motivates species-specific dormancy;
species-specific dormancy points to per-species flowering windows; flowering windows lead to seed dispersal.
There is no natural stopping point, and the discipline of holding scope to the written objectives is
difficult when each extension has a clear technical motive.

\medskip
A related lesson is that natural systems are much harder to parameterize convincingly than to implement
correctly. The L-system rewriting engine was straightforward to unit-test and verify. Choosing branch angles,
internode growth rates, dormancy thresholds, and seasonal colour transitions that together produce a
convincing plant took considerably longer and was guided almost entirely by visual iteration. Whilst
designing grammars to capture a specific plants structure became easier with practice and research, tuning
the individual parameters remained a time-consuming and challenging process.

\medskip
Choosing Rust for this project was beneficial for correctness --- the type system enforced important
invariants (for example, the \texttt{Plant} trait bounds and the generic L-system rule type) and caught
ownership errors at compile time that would otherwise have been runtime bugs in a garbage-collected language.
The main cost was the upfront design investment required to satisfy the borrow checker in the rendering and
scene controller code. With hindsight, a cleaner separation of mutable and immutable data early in the
architecture design would have reduced the number of refactors needed as the system grew.

% ──────────────────────────────────────────────────────────────────────────────────────────────────────────
\section{Legal, Social, Ethical and Professional Issues}

\subsection{Legal Issues}

All software dependencies used in this project are open-source under permissive licences, \texttt{cargo-deny}
was used to validate dependency licences.~\footnote{\url{https://github.com/EmbarkStudios/cargo-deny/}}. No
proprietary code, datasets, or assets are used. The project does not use, collect or store any personal or
private information. All third-party sources are properly cited throughout the report.

\medskip
The three plant grammars adapted from \textit{The Algorithmic Beauty of
Plants}~\cite{prusinkiewiczAlgorithmicBeautyPlants1990} are derived works. The adaptations are substantial:
each grammar has been re-parameterised, extended with seasonal behaviour, and re-implemented from scratch in
Rust rather than copied from any published source code. The original grammars are properly cited throughout
the report and in the grammar appendix. The FNV-1a hash function used in the voxel occupancy grid is a
well-known public domain algorithm; its use is acknowledged in the source code and
cited~\cite{fowlerFNVNonCryptographicHash2026}.

\subsection{Social Issues}

The primary social application of procedural vegetation systems is in entertainment (games) and simulation
(urban planning, landscape architecture, ecological modelling). In these contexts, high-quality vegetation
rendering can improve the realism of educational simulations and make ecological concepts more accessible to
non-specialist audiences. The system's open architecture (modular \texttt{Plant} trait) makes it
straightforward for domain experts to contribute new species grammars, which could serve as an educational
resource for students of both computer graphics and botany.

\medskip
A potential negative social impact is the use of realistic virtual environments as a substitute for
engagement with real natural environments. If virtual forests become more convenient to experience than
physical ones, this could reduce the public's motivation to protect real ecosystems. This concern is
speculative and difficult to quantify, however there has been some discussion over the importance of
vegetation in video games and its potential positive and negative impacts~\cite{truongGamersItGreen2018}.

\subsection{Ethical Issues}

The most significant ethical issue is energy consumption. GPU-intensive rendering consumes substantial
electrical power. At the scale of millions of simultaneously running games or simulations (a realistic figure
for popular game titles), the aggregate energy cost of procedural vegetation rendering has a non-trivial
impact. This is ethically notable because the aesthetic goal, realistic depiction of natural ecosystems, may
contribute to the degradation of real natural ecosystems through the energy demands of the computing
infrastructure required. The energy usage and efficiency of this project has not been quantified, efforts
have been taken to reduce the computational cost however the impact of these optimizations (LOD system,
culling) is unknown. If this project was to be deployed at scale the energy efficiency and potential
environmental impacts should be thoroughly assessed prior to deployment. The irony that realistic virtual
ecosystems may require energy production with real ecological costs (power station emissions, land use for
renewables) is an ethical tension that the games and simulation industry has not yet adequately
addressed~\cite{truongGamersItGreen2018,changGreenComputerVideo}.

\subsection{Professional Issues}

The project follows standard professional software engineering practices throughout: version control with Git
(providing a full audit trail of development decisions), reproducible builds via Nix (ensuring the software
can be rebuilt from source by any developer with access to the repository), and modular architecture with
documented interfaces.

\medskip
One professional limitation of this project as a solo piece of work is the absence of peer code review. All
design decisions, including several that were later refactored (the turtle's initial tight coupling to the
scene controller, the early absence of LOD), were made and reviewed by a single developer. In a team context
these issues would likely have been caught earlier.

% ──────────────────────────────────────────────────────────────────────────────────────────────────────────
\chapter{External Material}
\label{app:external}

\begin{itemize}
  \item \textbf{Leaf texture.} The leaf texture used is \textit{Leaf Set 024} from \textit{ambientCG.com},
    licensed under the Creative Commons CC0 1.0 Universal
    License.\footnote{\url{https://ambientcg.com/view?id=LeafSet024}}
  \item \textbf{Plant grammar designs.} The grammars for \textit{Capsella}, \textit{Lychnis}, and
    \textit{Mycelis} are adapted from grammars published in \textit{The Algorithmic Beauty of
    Plants}~\cite{prusinkiewiczAlgorithmicBeautyPlants1990}.
  \item \textbf{Ecosystem placement algorithm.} The kernel-based placement, self-thinning, and ecological
    succession algorithms are implementations of the methods described in Lane and
    Prusinkiewicz~\cite{laneGeneratingSpatialDistributions2002}.
  \item \textbf{Software libraries.} \texttt{wgpu}, \texttt{winit}, \texttt{egui}, \texttt{image},
    \texttt{glam}, \texttt{rand} and \texttt{bytemuck} are used as dependencies. No library source code has
    been copied into the project; all are used as external crates via the Cargo package manager.
  \item \textbf{Docs and examples.} The \texttt{wgpu} documentation and examples were consulted whilst
    implementing the rendering pipeline.~\footnote{\url{https://docs.rs/wgpu/27.0.1/wgpu/} and
    \url{https://github.com/gfx-rs/wgpu/tree/v27/examples#readme}}
\end{itemize}

No pre-written code, supervisor-provided materials, or external datasets were used. No human participant data
was collected.

% ──────────────────────────────────────────────────────────────────────────────────────────────────────────
\chapter{Turtle Actions}
\label{app:turtle}

\begin{table}[h]
  \centering
  \begin{tabular}{|l|l|l|}
    \hline
    \textbf{Action}       & \textbf{Effect}                               & \textbf{Notation} \\
    \hline
    \texttt{Branch(l, r)} & Draw cylinder: length $l$, base radius $r$    & $B(..)$, size parameters vary by plant \\
    \texttt{Leaf(w, h)}   & Draw leaf quad: width $w$, height $h$         & $L(..)$, size parameters vary by plant \\
    \texttt{Roll(a)}      & Rotate about heading axis by angle $a$ (rad)  & $/$ (clockwise), $\backslash$
    (anticlockwise) \\
    \texttt{Turn(a)}      & Rotate about normal axis by angle $a$ (rad)   & $+$ (right), $-$ (left) \\
    \texttt{Pitch(a)}     & Rotate about binormal axis by angle $a$ (rad) & $\wedge$ (up), $\mathbin{\&}$ (down) \\
    \texttt{Push}         & Save current turtle state                     & $[$ \\
    \texttt{Pop}          & Restore last saved state                      & $]$ \\
    \texttt{Cut}          & Terminate current branch                      & $\%$ \\
    \texttt{Colour(c)}    & Set current drawing colour to $c$             & $C(\mathbf{c})$ \\
    \texttt{Nop}          & No operation (used for non-geometric symbols) & \\
    \hline
  \end{tabular}
  \caption{Turtle graphics action set.}
  \label{tab:actions}
\end{table}

% ──────────────────────────────────────────────────────────────────────────────────────────────────────────
\chapter{Plant Reference}
\label{app:plants}

All figures in this section were generated using the default settings and a seed of 7, the ground and
background were set to black and a single plant was positioned at $(0, 0, 0)$. Images were then cropped and
stitched together. Unless otherwise specified grammars use the notation detailed in Table \ref{tab:actions}.

% ──────────────────────────────────────────────────────────────────────────────────────────────────────────
\section{Fern}
\label{app:fern}

\begin{figure}[h]
  \centering
  \resizebox{\textwidth}{!}{\includegraphics{fern}}
  \caption{Fern growth stages, each iteration is three steps apart.}
  \label{fig:fern}
\end{figure}

\textbf{Nonterminals:}
\begin{itemize}
  \item $\mathrm{L}(a, d, \delta)$ --- frond leaf node; age $a$, node distance $d$, branching depth $\delta$
  \item $\mathrm{Br}(a, \delta)$ --- branch segment; age $a$, branching depth $\delta$
  \item $\mathrm{Cy}(e, \delta)$ --- crown layer; expansion delay $e$, crown depth $\delta$
\end{itemize}

\textbf{Parameters:}
\begin{itemize}
  \item $\alpha$ --- branch angle (default $60^\circ$)
  \item $\theta_d$ --- frond droop angle (default $-10^\circ$)
  \item $a_{\max} = 7$ --- maximum leaf age
  \item $d_{\max} = 9$ --- maximum frond node distance
  \item $\delta_{\max} = 3$ --- maximum branching depth
  \item $a_b = 3$ --- minimum age for lateral branching
  \item $e_0 = 3$ --- crown expansion delay
  \item $\rho_c = \tfrac{2\pi}{3}$ --- crown radial separation angle
  \item $\mu = \tfrac{\pi}{6}$ --- crown pitch factor
\end{itemize}

\textbf{Grammar:}

Let $\Omega(\delta)$ denote the frond macro and $\varphi(d,\, \delta)$ denote the droop equation:
\begin{align*}
  \Omega(\delta)       &= \bigl[\, \wedge(\mu\sqrt{\delta})\; \mathrm{L}(a_b,\, \delta,\, 1)\, \bigr] \\
  \varphi(d,\, \delta) &= \theta_d\! \left(\tfrac{\delta_{\max} - \delta}{\delta_{\max}} + \tfrac{2d}{d_{\max}}\right)
\end{align*}

The grammar is defined as:
% tex-fmt: off
\begin{align*}
\omega \quad &: /\! \bigl(\tfrac{\pi}{2}\bigr)\; {-}\! \bigl(\tfrac{\pi}{2}\bigr)\; {\backslash}\! \bigl(\tfrac{\pi}{2}\bigr)\; \mathrm{Cy}(0,\, 1) \\
p_0    \quad &: \mathrm{L}(a, d, \delta) : a < a_{\max}\; \land\; \bigl(a < a_b\; \lor\; d > d_{\max}\; \lor\; \delta \ge \delta_{\max}\bigr)\; \rightarrow\;\mathrm{L}(a{+}1,\, d,\, \delta) \\
p_1    \quad &: \mathrm{L}(a, d, \delta) : a_b \le a \le a_{\max}\; \land\; d \le d_{\max}\; \land\; \delta < \delta_{\max} \\
             &  \qquad \; \rightarrow\;\mathbin{\&}(\varphi(d,\, \delta))\; \mathrm{Br}(1,\, \delta)\; \bigl[\, {-}(\alpha)\; \mathrm{L}(a_b{-}1,\, d,\, \delta{+}1)\, \bigr]\; \bigl[\, \mathrm{Br}(1,\, \delta)\; \mathrm{L}(a,\, d{+}1,\, \delta)\, \bigr]\; \bigl[\, {+}(\alpha)\; \mathrm{L}(a_b{-}1,\, d,\, \delta{+}1)\, \bigr] \\
p_2    \quad &: \mathrm{Br}(a, \delta) : a \le a_{\max}\; \rightarrow\;\mathrm{Br}(a{+}1,\, \delta) \\
p_3    \quad &: \mathrm{Cy}(e, \delta) : e > 0\; \rightarrow\;\mathrm{Cy}(e{-}1,\, \delta) \\
p_4    \quad &: \mathrm{Cy}(0, \delta)\; \rightarrow\;\Omega(\delta)\; {+}(\rho_c)\; \Omega(\delta)\; {+}(\rho_c)\; \Omega(\delta)\; {+}\!\left(\tfrac{\rho_c}{\delta{+}1}\right)\; \mathrm{Cy}(e_0,\, \delta{+}2)
\end{align*}
% tex-fmt: on

% ──────────────────────────────────────────────────────────────────────────────────────────────────────────
\section{Mint}
\label{app:mint}

\begin{figure}[h]
  \centering
  \resizebox{\textwidth}{!}{\includegraphics{mint}}
  \caption{Mint growth stages, each iteration is four steps apart.}
  \label{fig:mint}
\end{figure}

\textbf{Nonterminals:}
\begin{itemize}
  \item $A(t)$ --- main stem apex; node counter $t$
  \item $Q$ --- lateral sub-stem apex
  \item $I(l)$ --- internode segment; current length $l$
\end{itemize}

\textbf{Parameters:}
\begin{itemize}
  \item $\alpha_l$ --- leaf spread angle (default $70°$)
  \item $\alpha_b$ --- lateral sub-stem angle (default $35°$)
  \item $n_A$ --- maximum node count (default $6$)
  \item $\rho = \tfrac{\pi}{2}$ --- roll between internode pairs
\end{itemize}

\textbf{Grammar:}

Let $\Omega$ denote the leaf-node macro and let $\Lambda$ denote the lateral branch pair:
\begin{align*}
  \Omega  &=
  I\; [\, +(\alpha_l)\; L(w_l,h_l)\, ]\; [\, -(\alpha_l)\; L(w_l,h_l)\, ]\; /(\rho)\;
  I\; [\, +(\alpha_l)\; L(w_l,h_l)\, ]\; [\, -(\alpha_l)\; L(w_l,h_l)\, ]\; /(\rho) \\
  \Lambda &= [\, /(\pi)\; +(\alpha_b)\; Q\, ]\; [\, +(\alpha_b)\; Q\, ]
\end{align*}

The grammar is defined as:
% tex-fmt: off
\begin{align*}
\omega \quad &: A(0) \\
p_1    \quad &: A(t) : t < n_A,\; t \not\equiv 2 \; \rightarrow\; \Omega\; A(t{+}1) \\
p_2    \quad &: A(t) : t < n_A,\; t \equiv 2 \; \rightarrow\; \Omega\; \Lambda\; A(t{+}1) \\
p_3    \quad &: A(t) : t \ge n_A,\; t \not\equiv 2 \; \rightarrow\; \Omega \\
p_4    \quad &: A(t) : t \ge n_A,\; t \equiv 2 \; \rightarrow\; \Omega\; \Lambda \\
p_5    \quad &: Q \; \rightarrow\; \Omega_Q \quad \text{(same as $\Omega$ but with sub-stem $Q$)} \\
p_6    \quad &: I(l) \; \rightarrow\; I\! \left(\min(l \cdot r_i,\; i_{\max})\right)
\end{align*}
% tex-fmt: on

% ──────────────────────────────────────────────────────────────────────────────────────────────────────────
\section{Wildflower}
\label{app:wildflower}

\begin{figure}[h]
  \centering
  \resizebox{\textwidth}{!}{\includegraphics{wildflower}}
  \caption{Wildflower growth stages, each iteration is five steps apart.}
  \label{fig:wildflower}
\end{figure}

\textbf{Nonterminals:}
\begin{itemize}
  \item $P$ --- axiom
  \item $A(a)$ --- main stem apex; accumulated angle jitter $a$
  \item $I(a)$ --- lateral branch pair; branch half-angle $a$
  \item $Q$ --- branch tip
\end{itemize}

\textbf{Terminals:}
\begin{itemize}
  \item $F$ --- stem segment
\end{itemize}

\textbf{Parameters:}
\begin{itemize}
  \item $\alpha$ --- branch angle (default $40°$)
  \item $\phi \approx 137.5°$ --- golden divergence angle
  \item $\sigma_1 = 0.05,\; \sigma_2 = 0.50$ --- flower season window
  \item $\sigma_3 = 0.72$ --- seed dispersal season threshold
  \item $r_0 \sim \mathcal{U}(0,\, 10)$ --- initial apex jitter
  \item $j \sim \mathcal{U}(-a,\, a)$ --- per-step heading jitter (sampled using symbol parameter $a$)
  \item $r_b \sim \mathcal{U}(0°,\, \alpha)$ --- lateral branch angle
\end{itemize}

\textbf{Grammar:}
% tex-fmt: off
\begin{align*}
\omega \quad &: C_{\mathrm{stem}}\; P \\
p_1    \quad &: P \; \rightarrow\; [\, A(r_0)\, ] \\
p_2    \quad &: A(a) \; \rightarrow\; F\, F\, F\; I(40)\; {+}\!(j)\; /(\phi)\; A(a{+}1)\; [\, C_{\mathrm{top}}\, L(w_{\mathrm{top}},\, w_{\mathrm{top}})\, ]_{\sigma < \sigma_3} \\
p_3    \quad &: I(a) \; \rightarrow\; [\, {-}(a°)\; Q\, ]\; /(\phi)\; [\, {+}(a°)\; Q\, ] \\
p_4    \quad &: Q \xrightarrow{0.5} F\; Q\; [\, C_{\mathrm{leaf}}\, L(w_l,\, h_l)\, ] \\
p_5    \quad &: Q \xrightarrow{0.5} I(r_b)
\end{align*}
% tex-fmt: on

% ──────────────────────────────────────────────────────────────────────────────────────────────────────────
\section{\textit{Lychnis}}
\label{app:lychnis}

\begin{figure}[h]
  \centering
  \resizebox{\textwidth}{!}{\includegraphics{lychnis}}
  \caption{\textit{Lychnis} growth stages, each iteration is four steps apart.}
  \label{fig:lychnis}
\end{figure}

\textbf{Nonterminals:}
\begin{itemize}
  \item $\mathrm{A}(d)$ --- apex; delay counter $d$
  \item $\mathrm{I}(c)$ --- internode; remaining branch-segment pair count $c$
  \item $\mathrm{L}(s)$ --- leaf; growth step $s$
  \item $\mathrm{K}(s)$ --- flower bud; growth step $s$
\end{itemize}

\textbf{Terminals:}
\begin{itemize}
  \item $B$ --- branch segment
\end{itemize}

\textbf{Parameters:}
\begin{itemize}
  \item $\alpha_l = \tfrac{\pi}{3}$ --- leaf pitch
  \item $\alpha_b = \tfrac{\pi}{2}$ --- branch pitch
  \item $\rho = \tfrac{\pi}{2}$ --- roll between branch pairs
  \item $d_{\text{delay}} = 7$ --- apex expansion delay
  \item $s_{\max} = 10$ --- maximum growth steps
  \item $\sigma_1 = 0.10,\; \sigma_2 = 0.48$ --- bloom season window
\end{itemize}

\textbf{Grammar:}

Let $\Omega(d)$ denote the apex macro and let $\Lambda$ denote the leaf macro:
\begin{align*}
  \Omega(d) &= [\, \wedge(\alpha_b)\; \mathrm{A}(d)\, ] \\
  \Lambda   &= [\, \wedge(\alpha_l)\; \mathrm{L}(0)\, ]
\end{align*}

The grammar is defined as:
% tex-fmt: off
\begin{align*}
\omega \quad &: \mathrm{A}(d_{\text{delay}}) \\
p_0    \quad &: \mathrm{A}(d) : d < d_{\text{delay}}\; \rightarrow\;\mathrm{A}(d+1) \\
p_1    \quad &: \mathrm{A}(d) : d \ge d_{\text{delay}}\; \rightarrow\;B\; \mathrm{I}(20)\; \Lambda\; /(\rho)\; \Omega(0)\; /(\rho)\; \Lambda\; /(\rho)\; \Omega(d_{\text{delay}}{-}3)\; B\; \mathrm{I}(10)\; [\, \mathrm{K}(0)\, ]_{\sigma_1 < \sigma < \sigma_2} \\
p_2    \quad &: \mathrm{I}(c) : c > 0\; \rightarrow\;B\; B\; \mathrm{I}(c-1) \\
p_3    \quad &: \mathrm{L}(s) : s < s_{\max}\; \rightarrow\;\mathrm{L}(s+1) \\
p_4    \quad &: \mathrm{K}(s) : s < s_{\max}\; \rightarrow\;\mathrm{K}(s+1)
\end{align*}
% tex-fmt: on

% ──────────────────────────────────────────────────────────────────────────────────────────────────────────
\section{\textit{Mycelis}}
\label{app:mycelis}

\begin{figure}[h]
  \centering
  \resizebox{\textwidth}{!}{\includegraphics{mycelis}}
  \caption{\textit{Mycelis} growth stages, each iteration is five
  steps apart.}
  \label{fig:mycelis}
\end{figure}

\textbf{Nonterminals:}
\begin{itemize}
  \item $A(t)$ --- apex; delay counter $t$
  \item $I(t)$ --- internode; delay counter $t$
  \item $K(c)$ --- flower bud; countdown $c$
  \item $M$ --- internode marker
  \item $S$ --- basipetal signal
  \item $T$ --- signal propagation intermediate
  \item $W$ --- transient intermediate
  \item $G$ --- lateral bud
\end{itemize}

\textbf{Terminals:}
\begin{itemize}
  \item $F$ --- internode segment
  \item $V$ --- stable signal
  \item $L$ --- leaf
\end{itemize}

\textbf{Parameters:}
\begin{itemize}
  \item $\theta_b$ --- lateral branch angle (default $30°$); $\delta = \theta_b + \tfrac{\pi}{8}$
  \item $\phi \sim \mathcal{U}(\tfrac{\pi}{6},\, \pi)$ --- roll between internodes
  \item $c_{\max} = 7$ --- flower bud countdown
  \item $n_0 = 20$ --- internode delay
\end{itemize}

\textbf{Grammar:}
% tex-fmt: off
\begin{align*}
\omega    \quad &: I(n_0)\, F\, A(0) \\
p_1       \quad &: S < A(t)\; \rightarrow\;T\, V\, C_{\mathrm{flower}}\, K(c_{\max})\, C_{\mathrm{stem}} \\
p_2       \quad &: V < A(t)\; \rightarrow\;T\, V\, C_{\mathrm{flower}}\, K(c_{\max})\, C_{\mathrm{stem}} \\
p_3       \quad &: A(t): t > 0\; \rightarrow\;A(t-1) \\
p_4       \quad &: A(t)\; \rightarrow\;M\; [\, {+}(\delta)\; L\; {-}\!\left(\tfrac{\pi}{8}\right)\; G\, ]\; F\; /(\phi)\; A(2) \\
p_5       \quad &: S < M\; \rightarrow\;S \\
p_6       \quad &: S > T\; \rightarrow\;T \\
p_7       \quad &: T < G\; \rightarrow\;F\, A(2) \\
p_8       \quad &: V < M\; \rightarrow\;S \\
p_9       \quad &: T > V\; \rightarrow\;W \\
p_{10}    \quad &: W\; \rightarrow\;V \\
p_{11}    \quad &: I(t): t > 0\; \rightarrow\;I(t-1) \\
p_{12}    \quad &: I(t): t = 0\; \rightarrow\;S \\
p_{13}    \quad &: K(c): c > 0\; \rightarrow\;K(c-1) \\
\end{align*}
% tex-fmt: on

% ──────────────────────────────────────────────────────────────────────────────────────────────────────────
\section{\textit{Capsella}}
\label{app:capsella}

\begin{figure}[h]
  \centering
  \resizebox{\textwidth}{!}{\includegraphics{capsella}}
  \caption{\textit{Capsella} growth stages, each iteration is five
  steps apart.}
  \label{fig:capsella}
\end{figure}

\textbf{Nonterminals:}
\begin{itemize}
  \item $a(t)$ --- vegetative apex; steps remaining $t$
  \item $A$ --- flowering apex
  \item $u(t)$ --- lateral-branch angle accumulator; steps remaining $t$
  \item $I(l)$ --- internode; current length $l$
  \item $X(t)$ --- fruit pod; maturation countdown $t$
\end{itemize}

\textbf{Terminals:}
\begin{itemize}
  \item $K$ --- flower petal or fruit pod
\end{itemize}

\textbf{Parameters:}
\begin{itemize}
  \item $\alpha_l$ --- leaf pitch (default $70°$)
  \item $\alpha_b$ --- initial lateral-branch pitch (default $18°$)
  \item $\phi_g \approx 137.5°$ --- golden divergence angle
  \item $n_{\mathrm{veg}} = 15$ --- vegetative apex step count
  \item $n_u = 4$ --- angle accumulator step count
  \item $\tau = 9°$ --- pitch increment per accumulator step
  \item $t_x = 5$ --- pod maturation delay
  \item $\sigma_{\mathrm{pod}} = 0.35$ --- pod transition season threshold
\end{itemize}

\textbf{Grammar:}
% tex-fmt: off
\begin{align*}
\omega \quad &: I(i_{\mathrm{init}})\; a(n_{\mathrm{veg}}) \\
p_1    \quad &: a(t): t > 0\; \rightarrow\;[\, \mathbin{\&}(\alpha_l)\; L(w_l,h_l)\, ]\; /(\phi_g)\; I(i_{\mathrm{init}})\; a(t{-}1) \\
p_2    \quad &: a(t): t = 0\; \rightarrow\;[\, \mathbin{\&}(\alpha_l)\; L(w_l,h_l)\, ]\; /(\phi_g)\; I(i_{\mathrm{init}})\; A \\
p_3    \quad &: A\; \rightarrow\;[\, \mathbin{\&}(\alpha_b)\; u(n_u)\; I(i_{\mathrm{init}})\; I\! \left(\tfrac{i_{\mathrm{init}}}{2}\right)\; X(t_x)\; K\, K\, K\, K\, ]\; /(\phi_g)\; I(i_{\mathrm{init}})\; A \\
p_4    \quad &: I(l)\; \rightarrow\;I\! \left(\min(l \cdot r_i,\; i_{\max})\right) \\
p_5    \quad &: u(t): t > 0\; \rightarrow\;\mathbin{\&}(\tau)\; u(t{-}1) \\
p_6    \quad &: u(t): t = 0\; \rightarrow\;\mathbin{\&}(\tau) \\
p_7    \quad &: X(t): t > 0\; \rightarrow\;X(t{-}1)
\end{align*}
% tex-fmt: on

% ──────────────────────────────────────────────────────────────────────────────────────────────────────────
\section{Tree}
\label{app:tree}

\begin{figure}[h]
  \centering
  \resizebox{\textwidth}{!}{\includegraphics{tree}}
  \caption{Tree growth stages, each iteration is two steps apart.}
  \label{fig:tree}
\end{figure}

\textbf{Nonterminals:}
\begin{itemize}
  \item $T$ --- trunk apex
  \item $X$ --- branch apex
  \item $G$ --- leaf cluster
  \item $F(l,\, d)$ --- internode; length $l$, diameter $d$
\end{itemize}

\textbf{Parameters:}
\begin{itemize}
  \item $\beta$ --- branch angle (default $30.7°$)
  \item $f_{\mathrm{init}},\; f_{\mathrm{max}}$ --- initial and maximum internode length
  \item $d_{\mathrm{init}},\; d_{\mathrm{max}}$ --- initial and maximum internode diameter
  \item $r_f$ --- internode growth multiplier (default $1.35$)
  \item $f_{\mathrm{rand}}$ --- internode length jitter half-range (default $0.02$)
  \item $j_t \sim \mathcal{U}(-\varepsilon_t,\, \varepsilon_t)$ --- trunk roll jitter
  \item $j_b \sim \mathcal{U}(-\varepsilon_b,\, \varepsilon_b)$ --- branch roll jitter
  \item $p_{\mathrm{shed}}$ --- leaf shed probability
  \item $n_{\mathrm{shed}}$ --- minimum leaf age before shedding
\end{itemize}

\textbf{Grammar:}

Let $F_0 = F(f_{\mathrm{init}},\, d_{\mathrm{init}})$, $\Psi(\Theta) = C_{\mathrm{leaf}}\, [\, \Theta\; L^*\,
]\, C_{\mathrm{bark}}$ where $L^* \in \{L(w_l,h_l),\, \%\}$ with $\mathrm{P}(L^* = \%) = p_{\mathrm{shed}}
\cdot \mathbf{1}[\mathrm{age} \ge n_{\mathrm{shed}}]$. Let $\Phi(j)$ denote the branch macro:
\begin{align*}
  \Phi(j) =
  \bigl[\, \backslash(\beta)\; {-}(\beta)\; F_0\; X\; [{+}(2\beta)\; G]\, \bigr]\;
  \bigl[\, \backslash(4\beta)\; {-}(\beta)\; F_0\; X\; [{-}(2\beta)\; G]\, \bigr]\;
  \backslash(9\beta)\; {-}(\beta)\; F_0\; X\; [{+}(\beta)\; G]
\end{align*}

The grammar is defined as:
% tex-fmt: off
\begin{align*}
\omega \quad &: T \\
p_1    \quad &: T \; \rightarrow\; F_0\, F_0\, F_0\; /(j_t)\; \Phi(j_t) \\
p_2    \quad &: X \; \rightarrow\; F_0\; /(j_b)\; \Phi(j_b) \\
p_3    \quad &: G \; \rightarrow\; \Psi\! \bigl(\backslash(2\beta)\, {+}(2\beta)\bigr)\; \Psi\! \bigl(\backslash(2\beta)\, {+}(3\beta)\bigr)\; \Psi\! \bigl(/(\beta)\, \wedge(\beta)\, {-}(2\beta)\bigr)\; \Psi\! \bigl(/(2\beta)\, \mathbin{\&}(2\beta)\, {-}(2\beta)\bigr) \\
p_4    \quad &: F(l,\, d) \; \rightarrow\; F\! \left( \min(l \cdot r_f + r,\; f_{\mathrm{max}}),\; \min(d \cdot r_f,\; d_{\mathrm{max}}) \right) \quad r \sim \mathcal{U}({-}f_{\mathrm{rand}},\; f_{\mathrm{rand}})
\end{align*}
% tex-fmt: on

% ──────────────────────────────────────────────────────────────────────────────────────────────────────────
\section{Bush}
\label{app:bush}

\begin{figure}[h]
  \centering
  \resizebox{\textwidth}{!}{\includegraphics{bush}}
  \caption{Bush growth stages, each iteration is two steps apart.}
  \label{fig:bush}
\end{figure}

\textbf{Nonterminals:}
\begin{itemize}
  \item $T$ --- trunk apex
  \item $X$ --- first-level branch node
  \item $N$ --- recursive sub-branch node
\end{itemize}

\textbf{Terminals:}
\begin{itemize}
  \item $F$ --- internode segment
\end{itemize}

\textbf{Parameters:}
\begin{itemize}
  \item $\beta$ --- branch angle (default $35.7°$)
  \item $j_t \sim \mathcal{U}(-\varepsilon_t,\, \varepsilon_t)$ --- trunk roll jitter
  \item $j_b \sim \mathcal{U}(-\varepsilon_b,\, \varepsilon_b)$ --- branch roll jitter
  \item $s$ --- leaf scale
\end{itemize}

\textbf{Grammar:}

Let $\Lambda$ denote the leaf macro and $\Omega_{a,b,c,d}$ denote leaf configurations:
\begin{align*}
  \Lambda &= C_{\mathrm{leaf}}\, [\, L(w_l s,\, h_l s)\, ]\, C_{\mathrm{bark}} \\
  \Omega_a &= {+}(\beta)\; \Lambda \\
  \Omega_b &= {-}(\beta)\; \Lambda \\
  \Omega_c &= {+}(2\beta)\; \Lambda \\
  \Omega_d &= {-}(2\beta)\; \Lambda
\end{align*}

The grammar is defined as:
% tex-fmt: off
\begin{align*}
\omega \quad &: T \\
p_1    \quad &: T \; \rightarrow\; F\; /(j_t)\; X \\
p_2    \quad &: X \; \rightarrow\; F\; /(j_b)\; \bigl[\, \backslash(\beta)\; {-}(\beta)\; F\; N\; [\Omega_c]\, \bigr]\; \bigl[\, \backslash(4\beta)\; {-}(\beta)\; F\; N\; [\Omega_d]\, \bigr]\; \backslash(9\beta)\; {-}(\beta)\; F\; N\; [\Omega_a] \\
p_3    \quad &: N \; \rightarrow\; F\; F\; /(j_b)\; \bigl[\, \backslash(\beta)\; {-}(\beta)\; F\; N\; [\Omega_a\; \Omega_a]\, \bigr]\; \bigl[\, \backslash(4\beta)\; F\; [\Omega_b\; \Omega_b]\; N\; [\Omega_b\; \Omega_b]\, \bigr]\; \backslash(9\beta)\; {-}(\beta)\; F\; [\Omega_b\; \Omega_b]\; N
\end{align*}
% tex-fmt: on

% ──────────────────────────────────────────────────────────────────────────────────────────────────────────
\chapter{Figures}
\label{app:figures}

% ──────────────────────────────────────────────────────────────────────────────────────────────────────────
\section{GUI}
\label{app:gui}

\begin{figure}[h]
  \centering
  \resizebox{0.5\textwidth}{!}{\includegraphics{scene-gui}}
  \caption{Scene GUI.}
  \label{fig:scene-gui}
\end{figure}

\begin{figure}[h]
  \centering
  \resizebox{0.5\textwidth}{!}{\includegraphics{environment-gui}}
  \caption{Envrironment GUI.}
  \label{fig:environment-gui}
\end{figure}

\begin{figure}[h]
  \centering
  \resizebox{0.5\textwidth}{!}{\includegraphics{plants-gui}}
  \caption{Plants GUI.}
  \label{fig:plants-gui}
\end{figure}

\begin{figure}[h]
  \centering
  \resizebox{0.5\textwidth}{!}{\includegraphics{ecosystem-gui}}
  \caption{Ecosystem GUI.}
  \label{fig:ecosystem-gui}
\end{figure}

\begin{figure}[h]
  \centering
  \resizebox{0.5\textwidth}{!}{\includegraphics{camera-gui}}
  \caption{Camera GUI.}
  \label{fig:camera-gui}
\end{figure}

\begin{figure}[h]
  \centering
  \resizebox{0.5\textwidth}{!}{\includegraphics{display-gui}}
  \caption{Display GUI.}
  \label{fig:display-gui}
\end{figure}

\begin{figure}[h]
  \centering
  \resizebox{0.5\textwidth}{!}{\includegraphics{debug-hud}}
  \caption{Debug HUD.}
  \label{fig:debug-hud}
\end{figure}

\end{appendices}

% ──────────────────────────────────────────────────────────────────────────────────────────────────────────
\chapter{Unit Tests}
\label{app:tests}

The unit tests are evaluated using \texttt{cargo-nextest}.~\footnote{\url{https://nexte.st/}} Documentation
is available in the source code. The output of \texttt{cargo nextest run --cargo-profile release} is shown
below (the output has been slightly condensed for clarity, no tests were removed or altered).

\begin{verbatim}
PASS [0.002s] settings::tests::season_tint_is_continuous
PASS [0.002s] settings::tests::season_name_works_with_negative
PASS [0.002s] settings::tests::ecosystem_mark_dirty_increments_generation
PASS [0.002s] settings::tests::plant_type_growth_rate_positive
PASS [0.002s] settings::tests::season_name_wraps_at_1
PASS [0.002s] settings::tests::season_name_covers_all_quadrants
PASS [0.004s] settings::tests::plant_type_shade_tolerance_in_range
PASS [0.004s] settings::tests::environment_mark_dirty_increments_generation
PASS [0.004s] settings::tests::scene_mark_dirty_increments_generation
PASS [0.002s] settings::tests::season_tint_winter_is_cool
PASS [0.004s] settings::tests::season_tint_summer_is_white
PASS [0.002s] util::lsystem::tests::algae
PASS [0.002s] util::lsystem::tests::context_left_skips_branches
PASS [0.002s] util::lsystem::tests::context_left_skips_ignored_symbols
PASS [0.002s] util::lsystem::tests::context_left_only
PASS [0.002s] util::lsystem::tests::context_left_and_right_match
PASS [0.002s] util::lsystem::tests::context_no_match_wrong_left
PASS [0.002s] util::lsystem::tests::context_no_match_wrong_right
PASS [0.002s] util::lsystem::tests::context_right_nested_branch_skipped_after_mismatch
PASS [0.002s] util::lsystem::tests::context_right_only
PASS [0.002s] util::lsystem::tests::context_sensitive_has_priority_over_normal
PASS [0.002s] util::lsystem::tests::context_right_searches_multiple_branches
PASS [0.002s] util::lsystem::tests::context_sensitive_has_priority_over_parametric
PASS [0.003s] util::lsystem::tests::context_right_finds_symbol_in_branch
PASS [0.002s] util::lsystem::tests::evolve_zero_steps_unchanged
PASS [0.004s] util::lsystem::tests::context_no_match_at_left_edge
PASS [0.002s] util::lsystem::tests::nonterminal_without_rule_passes_through
PASS [0.002s] util::lsystem::tests::evolve_matches_repeated_step
PASS [0.004s] util::lsystem::tests::context_sensitive_has_priority_over_stochastic
PASS [0.002s] util::lsystem::tests::normal_fallback_when_context_unmatched
PASS [0.002s] util::lsystem::tests::parametric_ignored_symbol_does_not_satisfy_context
PASS [0.002s] util::lsystem::tests::parametric_reads_inline_value
PASS [0.003s] util::lsystem::tests::parametric_conditional_countdown
PASS [0.002s] util::lsystem::tests::stochastic_falls_back_to_normal
PASS [0.003s] util::lsystem::tests::parametric_context_matches_by_discriminant
PASS [0.003s] util::lsystem::tests::stochastic_never_fires_at_prob_0
PASS [0.002s] util::lsystem::tests::terminal_passes_through
PASS [0.004s] util::lsystem::tests::stochastic_always_fires_at_prob_1
PASS [0.002s] util::turtle::tests::branch_below_min_radius_skips_mesh
PASS [0.002s] util::turtle::tests::branch_generates_cylinder_geometry
PASS [0.002s] util::turtle::tests::combine_line_geometries_merges_indices
PASS [0.002s] util::turtle::tests::combine_mesh_geometries_offsets_indices
PASS [0.002s] util::turtle::tests::cut_at_depth_zero_is_ignored
PASS [0.002s] util::turtle::tests::cut_inside_branch_skips_to_pop
PASS [0.002s] util::turtle::tests::cut_only_affects_enclosing_branch
PASS [0.002s] util::turtle::tests::fresh_turtle_has_no_line_segments
PASS [0.003s] util::turtle::tests::branch_lod_zero_segments_skips_mesh
PASS [0.002s] util::turtle::tests::fresh_turtle_has_no_mesh
PASS [0.002s] util::turtle::tests::leaf_skip_prob_one_skips_all_leaves
PASS [0.002s] util::turtle::tests::leaf_generates_double_sided_quad
PASS [0.002s] util::turtle::tests::lod_segment_count_controls_branch_vertices
PASS [0.002s] util::turtle::tests::nested_push_pop_restores_to_outer_state
PASS [0.002s] util::turtle::tests::mesh_index_count_matches_geometry
PASS [0.002s] util::turtle::tests::multiple_leaves_accumulate
PASS [0.002s] util::turtle::tests::pitch_tilts_heading_upward
PASS [0.002s] util::turtle::tests::roll_does_not_change_heading
PASS [0.002s] util::turtle::tests::scale_multiplies_travel_distance
PASS [0.002s] util::turtle::tests::space_pruning_allows_well_separated_branches
PASS [0.002s] util::turtle::tests::space_pruning_considers_branch_width
PASS [0.002s] util::turtle::tests::space_pruning_inter_plant_branches_are_blocked
PASS [0.002s] util::turtle::tests::travel_produces_one_segment
PASS [0.002s] util::turtle::tests::space_pruning_disabled_allows_overlapping_branches
PASS [0.002s] util::turtle::tests::space_pruning_intra_plant_branches_never_block_each_other
PASS [0.003s] util::turtle::tests::push_pop_restores_position_for_travel
PASS [0.003s] util::turtle::tests::space_pruning_blocks_mid_segment_overlap
PASS [0.002s] util::turtle::tests::turn_rotates_heading_left
PASS [0.003s] world::ecosystem::tests::bell_is_decreasing
PASS [0.002s] world::ecosystem::tests::bell_at_one_is_near_zero
PASS [0.003s] world::ecosystem::tests::bell_at_zero_is_one
PASS [0.002s] world::ecosystem::tests::mixed_kappa_at_origin_is_negative_one
PASS [0.002s] world::ecosystem::tests::different_seeds_produce_different_placements
PASS [0.002s] world::ecosystem::tests::empty_species_returns_empty
PASS [0.002s] world::ecosystem::tests::mixed_kappa_boundary_switches_sign
PASS [0.002s] world::ecosystem::tests::self_thinning_non_overlapping_keeps_both
PASS [0.002s] world::ecosystem::tests::mixed_kappa_mid_range_is_positive
PASS [0.002s] world::ecosystem::tests::self_thinning_empty_stays_empty
PASS [0.003s] world::ecosystem::tests::mixed_kappa_inner_range_is_negative
PASS [0.003s] world::ecosystem::tests::mixed_kappa_far_range_near_zero
PASS [0.003s] world::ecosystem::tests::plants_placed_within_area
PASS [0.004s] world::ecosystem::tests::instances_have_zero_y_position
PASS [0.002s] world::ecosystem::tests::self_thinning_removes_smaller_overlapping_plant
PASS [0.005s] world::ecosystem::tests::generates_requested_plant_count
PASS [0.003s] world::ecosystem::tests::zero_num_plants_returns_empty
PASS [0.003s] world::ecosystem::tests::species_weights_respected
PASS [0.003s] world::ecosystem::tests::self_thinning_single_plant_unchanged
PASS [0.004s] world::ecosystem::tests::self_thinning_reduces_count
PASS [0.005s] world::ecosystem::tests::same_seed_produces_identical_results
PASS [0.002s] world::plants::tests::dormancy_positive_offset_shifts_onset_earlier
PASS [0.002s] world::plants::tests::dormancy_offset_wraps_across_year_boundary
PASS [0.003s] world::plants::tests::all_plants_produce_actions_at_iteration_0
PASS [0.003s] world::plants::tests::dormancy_max_in_full_winter
PASS [0.002s] world::plants::tests::lerp_at_max_age_returns_old
PASS [0.002s] world::plants::tests::lerp_at_age_zero_returns_young
PASS [0.002s] world::plants::tests::dormancy_smoothstep_at_recovery_midpoint
PASS [0.002s] world::plants::tests::lerp_at_half_age_is_midpoint
PASS [0.005s] world::plants::tests::all_plants_produce_actions_at_iteration_1
PASS [0.002s] world::plants::tests::lerp_clamps_past_max_age
PASS [0.002s] world::plants::tests::dormancy_wraps_at_season_one
PASS [0.002s] world::plants::tests::plant_instance_clone_preserves_fields
PASS [0.003s] world::plants::tests::dormancy_smoothstep_at_onset_midpoint
PASS [0.003s] world::plants::tests::dormancy_zero_before_onset
PASS [0.002s] world::plants::tests::plant_instance_default_values
PASS [0.002s] world::plants::tests::plant_instance_with_transform_preserves_delay
PASS [0.002s] world::plants::tests::plant_instance_with_transform_sets_fields
PASS [0.002s] world::plants::tests::plant_iteration_reported_correctly
PASS [0.002s] world::plants::tests::plant_max_iterations_positive
PASS [0.002s] world::plants::tests::plant_type_reported_correctly
PASS [0.003s] world::plants::tests::plant_properties_consistent
PASS [0.003s] world::plants::tests::set_iteration_to_same_value_does_not_change_actions

Summary [0.029s] 109 tests run: 109 passed, 0 skipped
\end{verbatim}
